\Needspace{6\baselineskip}
\begin{center}\hypertarget{native-266-article192}{}\pdfbookmark[2]{192.}{article-192}192.\end{center}
The method explained in the preceding article is chiefly suited to the
\emph{first} determination of an orbit that includes perturbations. As soon
as both the mean elliptic elements and the equations of the perturbations
are already known approximately, the most accurate determination, with
the aid of as many observations as possible, is most conveniently carried
out by the method of Art. 187, which needs no special explanation here.
If there are enough observations of the highest quality and they cover a
long interval of time, this method can in many cases also serve to determine
more accurately the masses of the perturbing planets, at least the larger
ones. If the mass assumed for a perturbing planet in calculating the
perturbations still seems insufficiently certain, introduce a further
unknown $\mu$ in addition to the six unknowns depending on the corrections
to the elements. Take the ratio of the corrected mass to the assumed mass
to be $1+\mu$ to $1$. We may then assume that the perturbations themselves
change in the same ratio. This evidently gives rise in each calculated
position to a new linear term that contains $\mu$ and can be worked out
without difficulty. Comparing the calculated positions with the observed
positions according to the \SourceLeaf{267}principles set out above will yield
the correction $\mu$ along with the corrections to the elements. Indeed, the
masses of \emph{several} planets can be determined more accurately in this
way, provided that they exert sufficiently large perturbations.
There is no doubt that the motions of the new planets, especially Pallas
and Juno, which undergo such large perturbations from Jupiter, will in a
few decades yield a highly accurate determination of Jupiter's mass by this
method. It may even be possible someday to determine the mass of one or
another of these new planets from the perturbations it exerts on
the others.
\begin{center}\rule{3cm}{.4pt}\end{center}
